% Auto-generated by export_petsc_thesis_outputs.py
\begin{table}[!htbp]
    \centering
    \small
    \setlength{\tabcolsep}{4pt}
    \resizebox{\textwidth}{!}{%
    \begin{tabular}{c | r r r | r r r | r}
        \toprule
        \multicolumn{1}{c|}{} & \multicolumn{3}{c|}{FETI-DP} & \multicolumn{3}{c|}{BDDC} & \multicolumn{1}{c}{Total run} \\
        Case & $n_{\lambda}$ & it & $\lambda_{\max}$ & $n_{\mathrm{phys}}$ & it & $\lambda_{\max}$ & time [s] \\
        \midrule
        A-S & 123 & 6 & 1.6 & 4160 & 6 & 1.5 & 2.162e-02 \\
        A-M & 251 & 7 & 2.0 & 16512 & 6 & 1.8 & 8.718e-02 \\
        A-L & 379 & 7 & 2.3 & 37056 & 7 & 2.0 & 4.400e-01 \\
        B-S & 309 & 9 & 1.7 & 8320 & 8 & 1.7 & 3.435e-02 \\
        B-M & 629 & 10 & 2.2 & 33024 & 9 & 2.0 & 1.509e-01 \\
        B-L & 949 & 11 & 2.5 & 74112 & 10 & 2.2 & 5.216e-01 \\
        C-S & 741 & 11 & 1.8 & 16512 & 10 & 1.7 & 4.171e-02 \\
        C-M & 1509 & 12 & 2.2 & 65792 & 10 & 2.0 & 2.340e-01 \\
        C-L & 2277 & 13 & 2.5 & 147840 & 11 & 2.2 & 7.747e-01 \\
        \bottomrule
    \end{tabular}%
    }
    \caption{Summary of PETSc benchmark results for FETI-DP and BDDC with the relative tolerance $10^{-10}$. The quantity $n_{\lambda}$ denotes the dimension of the FETI-DP multiplier problem, while $n_{\mathrm{phys}}$ denotes the number of active physical degrees of freedom in the BDDC branch. In all cases PETSc reported $\lambda_{\min}=1.0$, so the reported value $\lambda_{\max}$ is also the corresponding condition-number indicator. The time is the median total wall-clock time of the complete PETSc run over three repetitions.}
    \label{tab:ch6-petsc-main-results}
\end{table}
